Fix $k\ge 1$. For $N\ge 1$, a \emph{2-colouring} of $[N]:=\{1,2,\dots,N\}$ is a map $\chi:[N]\to\{0,1\}$. For a finite set $A\subseteq [N]$, define its nonempty subset-sum set
\[
\mathrm{FS}(A):=\Big\{\sum_{a\in S}a:\ \emptyset\ne S\subseteq A\Big\}\subseteq \mathbb N.
\]
We say $A$ is \emph{monochromatic-sum} for $\chi$ if
\begin{enumerate}
\item $\mathrm{FS}(A)\subseteq [N]$ (so every sum is assigned a colour), and
\item $\chi$ is constant on $\mathrm{FS}(A)$.
\end{enumerate}
Define $F(k)$ to be the minimal $N$ (if it exists) such that every 2-colouring of $[N]$ contains some $A\subseteq [N]$ with $|A|=k$ that is monochromatic-sum.

\emph{Ambiguity in the literal statement.} The original phrasing does not explicitly require $\mathrm{FS}(A)\subseteq [N]$, but since colours are only defined on $[N]$, the natural convention is to require it; otherwise the property is ill-posed. I adopt that convention.
